\subsection{Oscillatore armonico 3-dimensionale}
\[
V_i(x^{(i)}) = \frac{1}{2} m\omega^2{x^{(i)}}^2,
\]
\[
\mathcal{H} \ket{\psi_{n_1 n_2 n_3}} = E_{n_1 n_2 n_3} \ket{\psi_{n_1 n_2 n_3}}.
\]
Per tre dimensioni
\[
E_{n_i} = \hbar \omega_i \left( n_i + \frac{1}{2} \right),
\]
\[
E_{n_i} = \hbar \left( \omega_1 n_1 + \omega_2 n_2 + \omega_3 n_3 + \frac{1}{2} (\omega_1 + \omega_2 + \omega_3) \right).
\]
Nel caso di $ \omega_1 = \omega_2 = \omega_3 = \omega $, allora
\[
V(\vec{x}) = \frac{1}{2} m\omega^2 \sum_{i = 1}^3 {x^{(i)}}^2 = \frac{1}{2} m\omega^2 \left\Vert \vec{x} \right\Vert^2
\]
e ho potenziale armonico sferico, quindi il caso dell'oscillatore armonico isotropo (sferico).

In questo caso la degenerazione si ha da 
\[
E_{n_1 n_2 n_3} = \hbar \omega \left( n_ 1 + n_2 + n_3 + \frac{3}{2} \right) = \hbar \omega \left( N + \frac{3}{2} \right), \quad N = 0, 1, 2, \dots
\]
ed \`e calcolabile esplicitamente perch\'e per ogni $ N $ fissato scelgo
\[
n_1 = 0, 1, 2, \dots, N \text{ e poi ho su } n_2 \text{ la condizione} \\
\]
\[
n_2 = 0, 1, \dots, N - n_1, \text{ dopodich\'e} \\
\]
\[
n_3 \text{ \`e dato univocamente.}
\]
quindi degenerazione si ha quando
\[
\sum_{n_1 = 0}^N (N - n_1 + 1) = (N + 1) \sum_{n_1 = 0}^N 1 - \sum_{n_1 = 0}^N n_1 = \frac{1}{2} (N + 1)(N + 2).
\]